% STATUS: reviewed
\chapter{Architektūra}\label{ch:arch}

Šiame skyriuje sistema aprašoma kaip veikiančių dalių visuma: aštuoni konteineriai ir už ką
kiekvienas atsakingas, du tinklai su vieninteliu įėjimu ir vieninteliu išėjimu, keturios
taisyklės, kurios lemia, ką kuri dalis gali daryti, ir dvejos vidinės durys tarp API ir
asistento paslaugos. Toliau per šias dalis sekamos keturios užklausos --- pokalbio ėjimas,
ataskaita, programos dokumentas ir teisės aktų bazės atnaujinimas --- tiek, kiek reikia
suprasti, kas ką kviečia. Foninės užduotys čia aprašomos išsamiai, nes kiekviena ilga
sistemos operacija yra užduotis: eilė, užduočių vykdytojas, planuoklis ir kasdienė priežiūra.
Skyrius baigiamas technologijų rinkiniu, dviem compose failo režimais ir
architektūriniais sprendimais su jų priežastimis.

\section{Sistema ir jos konteineriai}\label{sec:arch:containers}

\product{} padeda atitikties pareigūnui nustatyti, kurie Lietuvos ir ES teisės aktai taikomi
finansų sektoriaus įmonei, išskiria tų aktų nustatomas pareigas, palygina jas su pačios
įmonės dokumentais ir parengia spragų ataskaitą. Iš tų pačių dokumentų jis rengia atitikties
programos dokumentų projektus, kuriuos žmogus redaguoja ir patvirtina. Naudotojas dirba
pokalbyje; struktūrizuoti puslapiai skirti peržiūrai ir administravimui. Visa tai veikia
viename serveryje kaip vienas Docker Compose projektas, aprašytas \file{docker-compose.yml}
faile: aštuoni konteineriai (\cref{tab:arch:containers}), iš kurių trys sukuriami iš pačios
saugyklos kodo --- API, asistento paslauga ir konsolė, --- o vienas, duomenų bazė, iš
standartinio atvaizdo su į jį sukompiliuotu plėtiniu.

\begin{table}[H]
\caption{Aštuoni \file{docker-compose.yml} konteineriai ir už ką kiekvienas atsakingas}\label{tab:arch:containers}
\begin{tabularx}{\textwidth}{@{}L{36mm}L{33mm}Y@{}}
\toprule
\hd{Konteineris} & \hd{Atvaizdas} & \hd{Už ką atsakingas} \\
\midrule
\code{legalops-endpoint} & \code{caddy:2.11-alpine} & Vienintelis įėjimas. Jis atveria serverio 80 prievadą ir kiekvieną užklausą pagal kelią nukreipia į konsolę, API, asistento paslaugą arba duomenų bazės įrankį; nustato saugumo antraštes ir rašo prieigos žurnalą JSON formatu. \\
\code{legalops-exit} & \code{vuknf/caddy-l4:2.11} & Vienintelis išėjimas: 443 prievado tarpinis serveris, kuris TLS ryšį persiunčia tik tada, kai šis nurodo vieną iš keturių leidžiamų adresų. \\
\code{legalops-vite} & sukuriamas iš \file{vite/} & Konsolė adresu \code{/console} ir viešoji svetainė šakniniu adresu: vienas Vite projektas su dviem programomis. Jame nėra nei duomenų, nei paslapčių. \\
\code{legalops-django} & sukuriamas iš \file{django/} & API (Django): paskyros ir sesijos, įmonės ir jų dokumentai, teisės aktų bazė, pareigų registras, ataskaitos, programos dokumentai, instrukcijos, pokalbių istorija, DI sąnaudos ir veiksmų žurnalas; įkelti failai diske; vidinės durys \code{/internal/assistant/*}. Šiame konteineryje veikia ir užduočių vykdytojo ciklas. \\
\code{legalops-worker} & \code{docker:27-cli} & Nieko savo. Tai Docker komandinės eilutės klientas, kuris paleidžia užduočių ciklą \code{legalops-django} konteineryje ir veikia tol, kol veikia tas ciklas. \\
\code{legalops-assistant} & sukuriamas iš \file{assistant/} & Asistento paslauga: kiekvienas modelio iškvietimas --- srautinis pokalbis su įrankiais, API užsakyti generavimai ir vektorizavimai, išorinių MCP serverių klientas. Jai priklauso fakulteto DI tinklų sąsajos (angl. \en{gateway}) raktas, bet ne duomenys. \\
\code{legalops-postgres} & sukuriamas iš \file{postgres/} & PostgreSQL su \code{pgvector} plėtiniu: visi duomenys, vektoriniai ir viso teksto indeksai. Jei duomenų katalogas tuščias, pirmą kartą paleidžiant įkeliamas \file{postgres/bootstrap.sql.gz}. \\
\code{legalops-dbgate} & \code{dbgate/dbgate:7.3.3-alpine} & Administravimo įrankis: duomenų bazės naršyklė adresu \code{/dbgate/}, pasiekiama tik administratoriams. \\
\bottomrule
\end{tabularx}
\end{table}

Kiekvieno konteinerio šakninė failų sistema yra tik skaitoma; tai, ką reikia rašyti, patenka
į \code{tmpfs} prijungimą arba į vieną iš dviejų serverio katalogų po \file{_DATA/} (duomenų
bazės failai ir įkelti failai). Įėjimas, išėjimas, API ir duomenų bazė veikia kaip uid~1000.
Užduočių vykdytojas ir duomenų bazės įrankis veikia kaip root, o konsolės ir asistento
paslaugos konteinerių \code{user} eilutė compose faile užkomentuota, todėl jie veikia kaip
numatytasis jų atvaizdo naudotojas. Compose failas nenustato nei išteklių ribų, nei būklės
patikrų, nei paleidimo tvarkos: paslaugos startuoja kartu, o įėjimas atsako 502, kol
paslauga, kuriai priklauso kelias, dar neveikia. Šių sprendimų saugumo pasekmės aptariamos
\cref{ch:security}; diegimas ir atnaujinimai --- \cref{ch:ops}.

\paragraph{API viduje.} Django projektas nenaudoja nė vienos Django kartu teikiamos
programos: nėra nei administravimo svetainės, nei \code{django.contrib.auth}, nei sesijų
karkaso, nei savų šablonų (\file{settings.py}). Jį sudaro vienuolika paties projekto programų
(\cref{tab:arch:apps}), dvi tarpinės programinės įrangos dalys --- \code{RequestGuardMiddleware}
(\file{common/middleware.py}), kuri atmeta NUL simbolį kelyje ar užklausos parametruose ir
reikšmę, kurios duomenų bazė negali priimti (išeinančią už leistinų ribų ar netaisyklingą),
paverčia atsakymu 400, ir Django \code{CommonMiddleware}, ---
bei viena visos paslaugos maršrutų lentelė \file{urls.py}. Kiekvienas atsakymas, taip pat ir
klaidos, yra JSON. Laiko reikšmės visur yra UTC su nurodyta laiko juosta. Patys maršrutai
išvardyti \cref{ch:api}, o lentelės --- \cref{ch:data}.

\begin{table}[H]
\caption{API programos, kaip jas įdiegia \file{settings.py}}\label{tab:arch:apps}
\begin{tabularx}{\textwidth}{@{}L{23mm}Y@{\hspace{5mm}}L{23mm}Y@{}}
\toprule
\hd{Programa} & \hd{Turinys} & \hd{Programa} & \hd{Turinys} \\
\midrule
\code{users} & Paskyros, sesijos, prisijungimas & \code{programme} & Programos dokumentai \\
\code{audit} & Veiksmų žurnalas & \code{chat} & Pokalbiai ir žinutės \\
\code{jobs} & Eilė ir užduočių vykdytojas & \code{prompts} & Instrukcijos ir jų versijos \\
\code{companies} & Įmonės ir jų dokumentai & \code{ops} & Būklės patikra, priežiūra, užklausų ribojimo langai \\
\code{corpus} & Teisės aktai, redakcijos, struktūros vienetai, pareigos & \code{assistant} & Asistento paslaugos klientas, DI sąnaudos, kainos, MCP serveriai \\
\code{analysis} & Ataskaitos ir veiksmų registras & & \\
\bottomrule
\end{tabularx}
\end{table}

\paragraph{Asistento paslaugos viduje.} Paslauga yra vienas Node procesas: Express programa
(\file{app.js}), kurią aptarnauja \file{server.js}. Jos katalogai atitinka tris jos darbus:
\file{routes/} laiko viešąjį pokalbio maršrutą ir dvi vidinių maršrutų grupes, \file{llm/} ---
tiekėjų registrą, pokalbio įrankius, srautą ir MCP klientą, o \file{services/} --- API
klientus ir nedidelius atmintyje laikomus registrus (kas prisijungęs, ėjimų limitas, vykstantys
ėjimai). Diske ji nieko nesaugo.

\section{Tinklai, įėjimas ir išėjimas}\label{sec:arch:network}

Konteineriai prijungti prie dviejų Docker tinklų, ir skirtumas tarp jų yra pagrindinė
diegimo saugumo savybė (\cref{fig:arch:containers}).

\begin{description}
  \item[Izoliuotas tinklas] \code{isolated} yra tiltinis tinklas, deklaruotas
    \texttt{internal: true}, su potinkliu 172.28.10.0/24. Tokiam tinklui Docker nesuteikia
    maršruto į išorę. Prie jo prijungtos visos paslaugos, išskyrus užduočių vykdytoją, ir
    gamybos režimu tai vienintelis tinklas, prie kurio prijungti API, asistento paslauga,
    duomenų bazė ir duomenų bazės įrankis.
  \item[Išorinis tinklas] \code{external} yra įprastas Docker tinklas, kurio compose failas
    pats nesukuria (jis deklaruotas \texttt{external: true}; jį sukuria
    \file{runUpdateThisStack.sh}, kai serveryje tinklo tokiu pavadinimu nėra). Gamybos režimu
    prie jo jungiasi tik įėjimas ir išėjimas.
    Įėjimas yra vienintelis konteineris su atvertu prievadu.
  \item[Be tinklo] Užduočių vykdytojas turi \texttt{network\_mode: none}. Jis bendrauja tik su
    į jį prijungtu Docker lizdu.
\end{description}

\begin{figure}[H]
\centering
\begin{tikzpicture}[x=1mm, y=1mm]
  \begin{scope}[on background layer]
    \draw[lane] (0,-14) rectangle (150,-86);
  \end{scope}
  \node[lanelabel, anchor=north west] at (1.5,-15.5) {izoliuotas tinklas --- vidinis, be maršruto į išorę};

  \node[ext] (browser) at (75,6) {naršyklė, per prieglobos TLS tarpinį serverį};
  \node[svc, text width=24mm] (endpoint) at (75,-14) {\textbf{endpoint}\\ Caddy, 80 prievadas};
  \node[svc, text width=22mm] (vite) at (16,-40) {\textbf{vite}\\ viešoji svetainė, konsolė};
  \node[svc, text width=22mm] (dbgate) at (134,-40) {\textbf{dbgate}\\ duomenų bazės įrankis};
  \node[svc, text width=24mm] (assistant) at (40,-63) {\textbf{assistant}\\ visi modelio iškvietimai};
  \node[svc, text width=24mm] (django) at (98,-63) {\textbf{django}\\ API ir užduočių vykdytojo ciklas};
  \node[store] (postgres) at (136,-63) {\textbf{postgres}};
  \node[svc, text width=22mm] (exit) at (62,-86) {\textbf{exit}\\ SNI leidžiamų adresų sąrašas};
  \node[svc, text width=22mm] (worker) at (136,-101) {\textbf{worker}\\ Docker CLI, be tinklo};

  \node[ext] (gw) at (13,-109) {ai.knf.vu.lt};
  \node[ext] (cellar) at (48,-109) {publications.europa.eu};
  \node[ext] (etar) at (83,-109) {www.e-tar.lt};
  \node[ext] (lb) at (106,-109) {www.lb.lt};

  \draw[flow] (browser) -- (endpoint);
  \draw[flow] (endpoint) -- node[lbl, pos=.55] {visa kita} (vite);
  \draw[flow] (endpoint) -- node[lbl, pos=.5] {\code{/api/assistant/*}} (assistant);
  \draw[flow] (endpoint) -- node[lbl, pos=.5] {\code{/api/*}} (django);
  \draw[flow] (endpoint) -- node[lbl, pos=.55] {\code{/dbgate/*}} (dbgate);

  \draw[flow] ([yshift=2.2mm]django.west) -- node[lbl, above] {\code{/internal/llm/*}} ([yshift=2.2mm]assistant.east);
  \draw[flow] ([yshift=-2.2mm]assistant.east) -- node[lbl, below, align=center] {\code{/api/*} naudotojo vardu,\\ \code{/internal/assistant/*}} ([yshift=-2.2mm]django.west);
  \draw[flow] (django) -- (postgres);
  \draw[flow] (dbgate) -- (postgres);

  \draw[flow] (assistant.south) -- (exit.north west);
  \draw[flow] (django.south west) -- (exit.north east);
  \draw[flow] (exit.south west) -- (gw.north);
  \draw[flow] (exit.south) -- (cellar.north);
  \draw[flow] (exit.south east) -- (etar.north);
  \draw[flow] (exit.east) -- (lb.north);

  \draw[opt] (worker.north west) -- node[lbl, pos=.5] {\texttt{docker exec}} (django.south east);

  \node[note, anchor=west] at ([xshift=1mm]endpoint.east) {abiejuose tinkluose;\\ vienintelis atvertas prievadas};
  \node[note, anchor=east] at ([xshift=-1mm]exit.west) {abiejuose tinkluose,\\ adresu 172.28.10.30};
\end{tikzpicture}
\caption{Konteineriai (jų pavadinimai be priešdėlio \code{legalops-}), izoliuotas tinklas ir
  kuri paslauga kurią gali pasiekti. Rodyklė eina nuo kviečiančiosios paslaugos į kviečiamąją.
  Užklausos patenka tik per įėjimą ir išeina tik per išėjimą, kuris persiunčia į keturis
  įvardytus adresus. Taškinė rodyklė nėra tinklo kelias: užduočių vykdytojas per Docker lizdą
  paleidžia užduočių ciklą API konteineryje.}\label{fig:arch:containers}
\end{figure}

\subsection{Vienintelis įėjimas}

Įėjimas yra Caddy atvirkštinis tarpinis serveris, sukonfigūruotas \file{endpoint/Caddyfile}.
Jis klausosi 80 prievado paprastu HTTP; TLS baigiasi prieglobos tarpiniame serveryje, esančiame
prieš serverį. Pirmiausia jis pliku 403 atsakymu atmeta užklausą, kurios antraštė \code{Origin}
nurodo kitą svetainę,
ir prie kiekvieno kito atsakymo prideda turinio saugumo politikos bei perdavimo saugumo antraštes
(\cref{ch:security}),
o tada nukreipia pagal kelią; laimi pirmas atitikmuo (\cref{tab:arch:ingress}).

\begin{table}[H]
\caption{Kaip įėjimas nukreipia užklausą, \file{endpoint/Caddyfile} tvarka}\label{tab:arch:ingress}
\begin{tabularx}{\textwidth}{@{}L{34mm}L{45mm}Y@{}}
\toprule
\hd{Kelias} & \hd{Kur nukreipiama} & \hd{Pastaba} \\
\midrule
\code{/dbgate/*} & \code{legalops-dbgate:3000} & Tik tada, kai API maršrutu \route{GET}{/api/auth/checkauth/admin} patvirtina, kad kviečiantysis yra prisijungęs administratorius (Caddy \code{forward_auth}). \\
\code{/api/assistant/*} & \code{legalops-assistant:3000} & Pokalbio srautas ir jo sustabdymo bei būsenos maršrutai. \\
\code{/api/*} & \code{legalops-django:8000} & Visa viešoji API. \\
visa kita & \code{legalops-vite:80} & Viešoji svetainė ir konsolė. \\
\bottomrule
\end{tabularx}
\end{table}

Iš šios lentelės išplaukia du dalykai. Asistento paslauga turi lygiai vieną viešą maršrutų
grupę --- pokalbį; visa kita, ką ji aptarnauja, yra vidiniai maršrutai. Be to, nė vienas
kelias, prasidedantis \code{/internal/}, nepersiunčiamas nei į API, nei į asistento paslaugą:
tokia užklausa nukrenta iki paskutinės eilutės, ir į ją atsako konsolės konteineris, kuris
aptarnauja tik viešąją svetainę ir konsolę.

\subsection{Vienintelis išėjimas}

Izoliuotame tinkle esantis konteineris negali atidaryti ryšio į internetą. Išėjimas yra
vienintelė leidžiama išimtis, ir ji siaura. Tai ketvirtojo lygmens tarpinis serveris
(\file{exit/Caddyfile}), kuris klausosi tik 443 prievado, perskaito serverio vardą, kurį TLS
klientas praneša atidarydamas ryšį (SNI lauką), ir, jei tas vardas yra jo sąraše, persiunčia
ryšį tuo pačiu vardu į 443 prievadą. Jis nieko neiššifruoja. Ryšys, nurodantis bet kurį kitą
adresą, neatitinka jokio maršruto ir nepersiunčiamas.

Paslauga leidžiamą adresą pasiekia per du įrašus, kurie abu turi egzistuoti: maršrutą
\file{exit/Caddyfile} faile ir \code{extra_hosts} eilutę naudojančios paslaugos aprašyme
\file{docker-compose.yml} faile, kuri adreso vardą susieja su fiksuotu išėjimo adresu
172.28.10.30. Tada paslauga jungiasi prie, jos manymu, paties adreso, o pasiekia išėjimą.
\Cref{tab:arch:egress} išvardija keturis adresus.

\begin{table}[H]
\caption{Išeinančio srauto leidžiamų adresų sąrašas: \file{exit/Caddyfile} adresai ir
  paslaugos, kurioms jie susieti}\label{tab:arch:egress}
\begin{tabularx}{\textwidth}{@{}L{38mm}L{30mm}Y@{}}
\toprule
\hd{Adresas} & \hd{Susietas paslaugai} & \hd{Kam naudojamas} \\
\midrule
\code{ai.knf.vu.lt} & asistento paslauga, API & DI tinklų sąsaja: kiekvienas modelio iškvietimas, kurį atlieka asistento paslauga. Compose failas šį vardą susieja ir API konteineriui, bet joks API kodas jo nekviečia. \\
\code{publications.europa.eu} & API & CELLAR, Leidinių biuro mašininė sąsaja: ES teisės aktai. \\
\code{www.e-tar.lt} & API & Teisės aktų registras (e-TAR). \\
\code{www.lb.lt} & API & Lietuvos banko puslapiai ir PDF failai. \\
\bottomrule
\end{tabularx}
\end{table}

Trys teisės aktų bazės adresai susieti API konteineriui, nes ten veikia užduočių vykdytojo
ciklas, kuris atsisiunčia teisės aktus. Atsisiuntimo modulis virš tinklo taisyklės prideda
savo taisyklę (\file{corpus/fetchers.py}): tik HTTPS per 443 prievadą, iš naujo tikrinama
kiekvieno peradresavimo metu, kad šaltinis negalėtų nukreipti užduočių vykdytojo į vieną iš
pačios sistemos paslaugų; kurie adresai iš viso pasiekiami, lemia leidžiamų adresų sąrašas.

Tokios tvarkos priežastis ta, kad dvi paslaugos, kuriose yra sistemos duomenys ir DI tinklų
sąsajos raktas, taip pat apdoroja ne jų parašytą tekstą --- įkeltus dokumentus, atsisiųstus
puslapius, modelio atsakymus, išorinių MCP serverių įrankius. Jei vieną iš jų priverstų
atidaryti užpuoliko pasirinktą ryšį, izoliuotas tinklas tam ryšiui neleidžia nueiti niekur,
išskyrus keturis adresus. Kaina sumokama eksploatacijoje: naujas teisės aktų šaltinis ar
išorinis MCP serveris veikia tik tada, kai jo adresas pridėtas abiejose vietose
(\cref{ch:ops}). EUR-Lex ir e-seimas sąmoningai nenaudojami, nes jie atmeta automatinius
klientus (\cref{ch:corpus}).

\section{Keturios projektavimo taisyklės}\label{sec:arch:rules}

Keturios taisyklės nustato darbo pasidalijimą tarp API ir asistento paslaugos. Kiekviena turi
priežastį ir kiekviena turi pasekmę, kurią galima patikrinti kode.

\begin{description}
  \item[Asistento paslaugai nepriklauso jokie duomenys.]
    Instrukcijos, pokalbiai, DI sąnaudų įrašai ir MCP serverių konfigūracijos saugomi API
    duomenų bazėje. Asistento paslauga neturi ryšio su duomenų baze --- jos konteineriui
    neperduodamas duomenų bazės adresas --- ir tai, ko reikia, skaito per API.
    \emph{Priežastis:} viena kiekvieno fakto kopija, saugoma pagal vieną prieigos taisyklių
    rinkinį, vienoje vietoje. \emph{Pasekmė:} paslaugą
    galima paleisti iš naujo ar pakeisti nieko neprarandant, išskyrus tuo metu vykstančius pokalbio
    ėjimus. Tai, ką ji laiko atmintyje --- ėjimų ribotuvo langai, prisijungusių kviečiančiųjų
    podėlis, vykstančių ėjimų registras, --- kuriama iš naujo nuo nulio, ir kode pažymėta, kad todėl
    paslauga sąmoningai veikia kaip vienas egzempliorius (\file{services/runs.js},
    \file{services/ratelimit.js}).

  \item[API niekada nekviečia modelio tiekėjo.]
    Užduočių vykdytojas ir žiniatinklio procesas prašo asistento paslaugos generavimo
    (\route{POST}{/internal/llm/generate}) arba vektorių
    (\route{POST}{/internal/llm/embed}); vienintelis klientas yra API modulis
    \file{assistant/gateway.py}. API išeinančias HTTP užklausas siunčia lygiai iš dviejų
    modulių: šio kliento ir teisės aktų atsisiuntimo modulio. \emph{Priežastis:} tiekėjai ir
    jų prisijungimo duomenys konfigūruojami vieną kartą, viename konteineryje.
    \emph{Pasekmė:} DI tinklų sąsajos raktas yra tik asistento paslaugos aplinkoje, modelis
    visur įvardijamas viena \code{provider:model} formos eilute, saugoma instrukcijos
    versijoje, o kiekvienas sistemos modelio iškvietimas eina per vieną kodo dalį, kuri gali
    matuoti jo trukmę, jį kartoti ir skaičiuoti jo žetonus (\cref{ch:models}).

  \item[Pokalbio įrankiai veikia per viešąją API su paties naudotojo prieigos raktu.]
    Kai pokalbio modelis iškviečia įrankį --- užpildyti įmonės profilį, perskaityti ataskaitą,
    pradėti analizę, --- asistento paslauga siunčia įprastą užklausą į \code{/api/*} su
    besikalbančio naudotojo prieigos raktu (\code{userClient} faile
    \file{services/django.js}). \emph{Priežastis:} taisyklė, kas ką gali matyti, --- naudotojas
    mato savo įrašus, administratorius visus įrašus, išskyrus kitų naudotojų pokalbius, ---
    vykdoma vieną kartą, API rodiniuose, ir įrankis negali jos apeiti. \emph{Pasekmė:}
    pokalbis negali padaryti nieko, ko jo naudotojas negalėtų padaryti konsolėje, kiekvienas
    jo atliktas pakeitimas užrašomas veiksmų žurnale to naudotojo vardu, o įrankis neturi
    maršruto nei į paskyras, nei į instrukcijas ar MCP nustatymus, nei jokio maršruto, kuris
    keistų teisės aktų bazę: pokalbis gali ieškoti ir skaityti teisės aktus bei pareigų
    registrą, bet niekada negali jų redaguoti (\cref{ch:chat}). Yra dvi sąmoningos išimtys, abi
    per toliau aprašytas vidines duris: paties asistento atsakymai ir pavadinimas, kurį
    pokalbiui suteikia modelis.

  \item[Kiekvienas modelio iškvietimas registruojamas.]
    Iškvietimas palieka eilutę lentelėje \code{llm_calls} su modeliu, žetonų skaičiais, kaina
    JAV doleriais pagal DI tinklų sąsajos kainoraštį, instrukcijos versija ir rezultatu. API
    eilutę įrašo pati po kiekvieno savo užsakyto generavimo ir kiekvieno vektorizavimo paketo
    (\code{record_call} faile \file{assistant/usage.py}); asistento paslauga kiekvieną
    pokalbio ėjimą kartu su jo stenograma praneša per vidines duris. \emph{Priežastis:} kaina
    ir atskaitomybė --- biudžeto laikymąsi galima užtikrinti, o atsakymą atsekti tik turint išsamų
    įrašą. \emph{Pasekmė:} naudotojo dienos biudžetas, DI sąnaudų puslapis ir kiekvieno
    užduoties etapo kaina skaičiuojami iš šios vienos lentelės. Reikia įvardyti dvi taisyklės
    ribas: paieškos užklausos vektorizavimas (keli žetonai) neregistruojamas, o nepavykęs
    generavimas iškelia klaidą prieš įrašant eilutę, todėl užduočių etapų eilutės yra tik tų
    iškvietimų, kurie sulaukė atsakymo.
\end{description}

\section{Dvejos vidinės durys}\label{sec:arch:doors}

API ir asistento paslauga kviečia viena kitą, ir nė vienas iš šių kvietimų neatviras niekam
kitam. Kiekviena pusė turi \term{vidines duris}: maršrutų grupę po \code{/internal/}, kuri
egzistuoja tik kitai paslaugai (\cref{tab:arch:doors}).

\begin{xltabular}{\textwidth}{@{}L{62mm}Y@{}}
\caption{Kas eina pro kiekvienas vidines duris}\label{tab:arch:doors}\\
\toprule
\hd{Maršrutas} & \hd{Kas eina pro duris} \\
\midrule
\endfirsthead
\toprule
\hd{Maršrutas} & \hd{Kas eina pro duris} \\
\midrule
\endhead
\bottomrule
\endlastfoot
\multicolumn{2}{@{}L{\textwidth}@{}}{\hd{API durys} --- kviečia asistento paslauga (\code{internalFetch} faile \file{services/django.js}), aptarnauja \file{assistant/api/internal_views.py}} \\
\route{GET}{/internal/assistant/prompt} & Aktyvi vienos instrukcijos versija su tekstu ir modeliu. Jei instrukcija aktyvios versijos neturi, atsakoma 404, o pokalbis atsako 503, užuot veikęs pagal tekstą, kurio niekas negali matyti. \\
\route{GET}{/internal/assistant/budget} & Kiek naudotojas šiandien dar gali išleisti; klausiama prieš kiekvieną pokalbio ėjimą. \\
\route{POST}{/internal/assistant/turn-log} & Viena \code{llm_calls} eilutė pokalbio ėjimui: žetonai, iškviesti įrankiai, rezultatas ir stenograma. \\
\route{POST}{/internal/assistant/threads/{id}/messages} & Asistento atsakymas --- tekstas, įrankių iškvietimai ir įrankių rezultatai, --- įrašomas ėjimui pasibaigus. Tik pro šias duris galima įrašyti \code{assistant} rolės žinutę. \\
\route{POST}{/internal/assistant/threads/{id}/title} & Trumpas pavadinimas, kurį modelis parašė naujam pokalbiui; pritaikomas tik tol, kol pokalbis dar turi automatinį pavadinimą. \\
\route{GET}{/internal/assistant/mcp-servers} & Įjungti MCP serveriai su antraščių reikšmėmis, kuriomis prie jų tapatinamasi, --- vienintelė vieta, kur tos reikšmės palieka API. \\
\midrule
\multicolumn{2}{@{}L{\textwidth}@{}}{\hd{Asistento paslaugos durys} --- kviečia API (\file{assistant/gateway.py}), aptarnauja \file{routes/internal.js} ir \file{routes/internalMcp.js}} \\
\route{POST}{/internal/llm/generate} & Vienas generavimas: instrukcija, pasirinktiniai nurodymai ir bendri duomenys, modelis ir, jei reikia, JSON schema, kurią atsakymas turi atitikti. Grąžina tekstą arba objektą kartu su sąnaudomis. \\
\route{POST}{/internal/llm/embed} & Ne daugiau kaip 256 tekstų paketas (\code{MAX_EMBED_TEXTS}) vektorizuoti. \\
\route{GET}{/internal/llm/models} & Sukonfigūruoti tiekėjai ir DI tinklų sąsajos modeliai su kainomis. \\
\route{POST}{/internal/mcp/list-tools} & Prisijungti prie vieno MCP serverio ir išvardyti jo įrankius: administratoriaus ryšio patikra. \\
\end{xltabular}

Pirmoji trečiosios projektavimo taisyklės išimtis matyti lentelėje. Viešoji API leidžia
naršyklei įrašyti tik \code{user} rolės žinutes (\code{GUI_WRITABLE_ROLES} faile
\file{chat/api/views.py}); asistento žinutę, kurios įrankių rezultatus modelis vėl mato kitame
ėjime, galima įrašyti tik pro vidines duris. Todėl naršyklė niekada negali į pokalbį įterpti
suklastoto įrankio rezultato.

Duris saugo trys dalykai. Jos pasiekiamos tik izoliuotame tinkle. Įėjimas nepersiunčia nė
vieno \code{/internal/} kelio (\cref{sec:arch:network}). Ir kiekviena užklausa turi pateikti
bendrą paslaptį.

\begin{codebox}{Kaip tikrinama bendra paslaptis}
Paslaptis yra viena reikšmė, \code{ASSISTANT_INTERNAL_SECRET}, kurią atnaujinimo scenarijus
sugeneruoja kaip atsitiktinę šešioliktainę eilutę, o compose failas perduoda abiem
konteineriams. Kviečiantysis ją siunčia antraštėje \code{X-Internal-Secret}.

\begin{itemize}
  \item \textbf{API pusėje} dekoratorius \code{require_internal}
    (\file{assistant/api/internal_views.py}) apgaubia kiekvieną durų rodinį. Jis pateiktą ir
    laukiamą reikšmes palygina kaip baitus funkcija \code{hmac.compare_digest}, kuri užtrunka
    vienodai, kad ir kur reikšmės skirtųsi, todėl atsakymo trukmė spėliotojui nieko nepasako
    apie tai, kokia spėjimo dalis buvo teisinga.
  \item \textbf{Asistento paslaugoje} tarpinė programinė įranga \code{requireInternal}
    (\file{middleware/internal.js}) stovi prieš kiekvieną \code{/internal} kelią, dar prieš
    išnagrinėjant užklausos turinį. Ji lygina funkcija \code{timingSafeEqual}; kito ilgio
    reikšmė atmetama prieš palyginimą, o tai atskleidžia ilgį ir nieko daugiau.
  \item \textbf{Abi pusės klaidos atveju uždaro prieigą.} Kai sukonfigūruota paslaptis tuščia,
    atmetamas kiekvienas vidinis kvietimas; nė viena pusė nelaiko, kad \enquote{paslaptis
    nenustatyta} reiškia \enquote{paslapties nereikia}. Asistento paslauga tai dar ir
    praneša savo žurnale paleidžiama.
\end{itemize}
Į atmestą kvietimą atsakoma 403 su kodu \code{INTERNAL_ONLY}.
\end{codebox}

\section{Keturios užklausos, sekamos per sistemą}\label{sec:arch:flows}

Toliau aprašyti keturi srautai yra tie, kuriais remiasi visi kiti skyriai. Jie nubraižyti kaip
sekos: laikas eina žemyn, ištisinė rodyklė yra pagrindinio kelio žingsnis, punktyrinė rodyklė
--- užklausa kartu su atsakymu. Kas vyksta kiekvieno žingsnio viduje --- ką modelis gauna ir
kas daroma su jo atsakymu, --- aptariama kiekviename poskyryje nurodytuose skyriuose.

\subsection{Pokalbio ėjimas}

Pokalbio ėjimas yra vienintelis srautas, kuriame naršyklė bendrauja su asistento paslauga
(\cref{fig:arch:chat}). Konsolė siunčia \route{POST}{/api/assistant/chat} su naudotojo
prieigos raktu ir tik nauja žinute: patį pokalbį saugo API, todėl istorija, kurią mato
modelis, visada yra išsaugotoji, o ne tokia, kokią ją pateikia naršyklė.

\begin{figure}[H]
\centering
\begin{tikzpicture}[x=1mm, y=1mm]
  \begin{scope}[on background layer]
    \draw[lane] (53,-76) rectangle (150,-106);
  \end{scope}
  \node[lanelabel, anchor=south west] at (60,-106) {kartojama, kol modelis kviečia įrankius};

  \node[svc] (hc) at (10,0) {konsolė};
  \node[svc] (ha) at (58,0) {asistento paslauga};
  \node[svc] (hd) at (108,0) {API (Django)};
  \node[ext] (hg) at (140,0) {DI tinklų sąsaja};
  \foreach \x in {10,58,108} \draw[draw=rule, line width=.6pt] (\x,-4) -- (\x,-131);
  \draw[draw=rule, line width=.6pt] (140,-4) -- (140,-78);

  \draw[flow] (10,-13) -- node[lbl, above, align=center] {\code{POST}~\code{/api/assistant/chat}\\ raktas, nauja žinutė} (58,-13);
  \draw[call] (58,-21) -- node[lbl, above] {\code{/api/auth/me}: kas klausia} (108,-21);
  \node[chk, minimum height=5mm, inner ysep=2pt] at (58,-28.5) {ėjimų limitas, atmintyje};
  \draw[call] (58,-37) -- node[lbl, above] {\code{/internal/assistant/budget}} (108,-37);
  \draw[call] (58,-44) -- node[lbl, above] {\code{/internal/assistant/prompt}} (108,-44);
  \draw[call] (58,-54) -- node[lbl, above, align=center] {\code{/api/chat/threads/…}: pokalbis, jo\\ istorija; įrašoma nauja žinutė} (108,-54);
  \draw[call] (58,-64) -- node[lbl, above, align=center] {\code{/api/companies/…}, \code{/api/reports}:\\ įmonė ir jos ataskaitos} (108,-64);
  \draw[call] (58,-71) -- node[lbl, above] {\code{/internal/assistant/mcp-servers}} (108,-71);

  \node[ai, text width=15mm, minimum height=24mm] (m) at (140,-91) {modelis atsako per ne daugiau kaip 24 žingsnius};
  \draw[flow] (58,-82) -- node[lbl, above, pos=.3] {nurodymai, istorija, įrankiai} (m.west |- 0,-82);
  \draw[flow] (m.west |- 0,-89) -- node[lbl, above, pos=.72] {įrankio iškvietimas} (58,-89);
  \draw[call] (58,-94) -- node[lbl, above] {\code{/api/*} su naudotojo raktu} (108,-94);
  \draw[flow] (58,-99) -- node[lbl, above, pos=.3] {įrankio rezultatas} (m.west |- 0,-99);

  \draw[flow] (58,-112) -- node[lbl, above] {atsakymas įvykių srautu} (10,-112);
  \draw[call] (58,-120) -- node[lbl, above] {\code{…/threads/{id}/messages}: atsakymas} (108,-120);
  \draw[call] (58,-127) -- node[lbl, above] {\code{/internal/assistant/turn-log}} (108,-127);
\end{tikzpicture}
\actorlegend
\caption{Pokalbio ėjimas. Asistento paslauga nustato kviečiantįjį ir naudotojo vardu per viešąją
  API įkelia kontekstą, pro vidines duris perskaito instrukciją, biudžetą ir MCP serverius ir
  paleidžia modelį su įrankiais. Kiekvienas įrankis yra užklausa viešajai API su naudotojo
  prieigos raktu. Atsakymas ir sąnaudų įrašas grįžta pro vidines duris.}\label{fig:arch:chat}
\end{figure}

Maršrutas (\file{routes/chat.js}) prieš iškviesdamas bet kurį modelį praleidžia ėjimą per
savo patikras. Kviečiančiojo prieigos raktas tikrinamas per \route{GET}{/api/auth/me}, o
atsakymas 15~sekundžių laikomas podėlyje (\code{IDENTITY_CACHE_MS}); kviečiantysis be
galiojančio rakto atmetamas. Ėjimų ribotuvas leidžia 12 ėjimų vienam naudotojui per minutę
(\code{RATE_LIMIT_TURNS}). Naudotojo dienos biudžeto klausiama API, ir išnaudotas biudžetas
ėjimą atmeta; jei nepavyksta pati užklausa, ėjimas praleidžiamas, kad nebūtų užblokuoti visi
pokalbiai. Aktyvi instrukcijos \prompt{chat.assistant} versija skaitoma iš API, o jos nesant
pokalbis atsako 503: tada asistentas išjungtas sąmoningai.

Tada paslauga naudotojo vardu įkelia pokalbį, jo istorijos šaką, įmonę, apie kurią pokalbis,
ir tos įmonės ataskaitas, įrašo naudotojo žinutę ir atveria įjungtų MCP serverių įrankius. Ji
sudaro nurodymus ir iškviečia modelį DI tinklų sąsajoje su pridėtais įrankiais. Modelis gali
atsakyti arba iškviesti įrankį; įrankio iškvietimas virsta užklausa viešajai API naudotojo
prieigos raktu, jo rezultatas grąžinamas modeliui, ir tai kartojama ne daugiau kaip 24 modelio
žingsnius (\code{MAX_STEPS}), iš kurių paskutinis vykdomas be įrankių, kad ėjimas baigtųsi
tekstu. Atsakymas naršyklei siunčiamas srautu, kol yra rašomas.

Ėjimui pasibaigus paslauga įrašo atsakymą ir užregistruoja ėjimą. Ėjimas gyvuoja ilgiau nei
jį pradėjusi užklausa: jei naršyklė perkraunama ar skirtukas uždaromas, modelis tęsia darbą ir
atsakymas įrašomas; be paties ėjimo laiko ribų, anksčiau laiko jį užbaigti gali tik
sustabdymo mygtukas, naujesnis ėjimas tame pačiame pokalbyje arba išjungimas. Pirmajame
pokalbio ėjime, kol pokalbis dar turi automatinį pavadinimą, antras, nedidelis modelio
iškvietimas suteikia pokalbiui pavadinimą. Įrankiai, instrukcija ir abu modelio iškvietimai
aptariami \cref{ch:chat}.

\subsection{Ataskaita}

Ataskaita gerokai per ilga vienai užklausai: tai užduotis (\cref{fig:arch:report}).
\route{POST}{/api/reports} patikrina, ar kviečiantysis gali matyti įmonę, ar teisės aktų bazė
netuščia, ar nevyksta kita tos įmonės ataskaita ir ar naudotojui liko biudžeto; tada sukuria
ataskaitos eilutę būsenoje \code{queued}, į eilę įtraukia \code{report.analyze} rūšies užduotį
su ataskaitos id kaip turiniu ir iškart atsako.

\begin{figure}[H]
\centering
\begin{tikzpicture}[x=1mm, y=1mm,
    cell/.style={minimum height=5.6mm, inner ysep=1.5pt, inner xsep=2pt, text width=18mm}]
  \node[svc, text width=17mm] (hc) at (10,0) {konsolė ar pokalbio įrankis};
  \node[svc] (hd) at (44,0) {API (Django)};
  \node[store] (hj) at (70,0) {\code{jobs}};
  \node[svc] (hw) at (100,0) {vykdytojo ciklas};
  \node[svc, text width=17mm] (ha) at (140,0) {asistento paslauga};
  \foreach \x in {44,70,100} \draw[draw=rule, line width=.6pt] (\x,-5) -- (\x,-128);
  \draw[draw=rule, line width=.6pt] (10,-6) -- (10,-128);
  \draw[draw=rule, line width=.6pt] (140,-6) -- (140,-128);

  \draw[flow] (10,-14) -- node[lbl, above] {\code{POST}~\code{/api/reports}} (44,-14);
  \draw[flow] (44,-21) -- node[lbl, above] {užduotis eilėje} (70,-21);
  \draw[call] (100,-28) -- node[lbl, above] {paima užduotį} (70,-28);

  \node[svc, cell] at (89.5,-36) {profile};
  \node[ai, cell] at (110.5,-36) {documents};
  \node[ai, cell] at (89.5,-43) {citations};
  \node[ai, cell] (p1) at (110.5,-43) {applicability};
  \draw[call] (121.2,-39.5) -- node[lbl, above] {\code{/internal/llm/*}} (140,-39.5);

  \draw[flow] (100,-53) -- node[lbl, above, align=center] {užduotis atlikta; ataskaita\\ laukia būsenoje \code{scope_review}} (70,-53);
  \draw[call] (10,-61) -- node[lbl, above] {tikrina \code{GET}~\code{/api/reports/{id}}} (44,-61);

  \node[man, text width=17mm] (person) at (10,-73) {žmogus patvirtina apimtį};
  \draw[flow] (person.east) -- node[lbl, above] {\code{POST}~\code{…/scope}} (44,-73);
  \draw[flow] (44,-80) -- node[lbl, above] {vėl eilėje} (70,-80);
  \draw[call] (100,-87) -- node[lbl, above] {paima užduotį} (70,-87);

  \node[ai, cell] at (89.5,-95) {documents};
  \node[ai, cell] at (110.5,-95) {obligations};
  \node[ai, cell] at (89.5,-102) {gap\_analysis};
  \node[ai, cell] at (110.5,-102) {internal};
  \node[ai, cell] at (89.5,-109) {action\_plan};
  \node[ai, cell] at (110.5,-109) {summary};
  \draw[call] (121.2,-102) -- node[lbl, above] {\code{/internal/llm/*}} (140,-102);

  \draw[flow] (100,-118) -- node[lbl, above] {užduotis atlikta; ataskaita baigta} (70,-118);
  \draw[call] (10,-125) -- node[lbl, above] {skaito, eksportuoja} (44,-125);
\end{tikzpicture}
\actorlegend
\caption{Ataskaita, pradėta su apimties peržiūra. Užduotis vykdoma du kartus: pirmasis
  vykdymas baigiasi, kai teisės aktai gauna sprendimus, ir palieka ataskaitą laukti; žmogaus
  patvirtinimas vėl įtraukia į eilę tos pačios rūšies užduotį, o antrasis vykdymas vertina
  pareigas. Langeliai užduočių vykdytojo linijoje yra užduoties etapai; geltonas etapas kviečia
  modelius per asistento paslaugą. Antrasis vykdymas iš naujo perskaito įmonę, antrą kartą
  užfiksuoja profilį ir vėl atlieka dokumentų etapą, kuris
  inventorius randa podėlyje; etapas \code{internal} vykdomas tik ataskaitai, pradėtai su
  įmonės vidaus reikalavimais. Be peržiūros abi dalys atliekamos per vieną
  vykdymą.}\label{fig:arch:report}
\end{figure}

Užduočių vykdytojas paima užduotį ir paleidžia jos doroklę (angl. \en{handler}) ---
\code{analyze} faile \file{analysis/pipeline.py}. Ataskaita rengiama dviem dalimis. Pirmoji
nusprendžia, kurie teisės aktai taikomi, ir yra pigi: etapai \code{profile},
\code{documents}, \code{citations} ir \code{applicability}. Antroji vertina kiekvieną tų aktų
pareigą, ir čia išleidžiama daugiausia pinigų: \code{obligations}, \code{gap_analysis},
\code{action_plan} ir \code{summary}, o tarp spragų analizės ir plano --- etapas
\code{internal}, jei ataskaita pradėta su įmonės vidaus reikalavimais. Kiekvienas etapo
modelio iškvietimas keliauja iš užduočių vykdytojo pro asistento paslaugos duris, o iš ten ---
į DI tinklų sąsają.

\begin{humanbox}{Sustojimas tarp dviejų dalių}
Ataskaita, pradėta su peržiūros parinktimi, savaime netęsiama. Kai pirmoji dalis baigiasi,
doroklė nustato ataskaitai būseną \code{scope_review} ir grįžta; užduotis baigiasi kaip
atlikta, ir niekas nebevykdoma. Konsolė, kuri, kol ataskaita aktyvi, kas tris sekundes
tikrina \route{GET}{/api/reports/{id}}, rodo teisės aktus ir ką ataskaita su kiekvienu
darytų. Žmogus gali pakeisti atranką ir pamatyti jos kainą (\route{POST}{/api/reports/{id}/scope}
be \code{approved}), o tada ją patvirtina (tas pats maršrutas su \code{approved}).
Patvirtinimas pažymimas ataskaitoje su asmeniu ir laiku, įrašomas į veiksmų žurnalą su
kiekvienu padarytu pakeitimu, ir į eilę įtraukiama antra \code{report.analyze} užduotis.
Doroklė antrąjį vykdymą nuo pirmojo atskiria pagal šią žymą. Be peržiūros parinkties doroklė
laiką pažymi pati, be asmens, ir iškart tęsia.
\end{humanbox}

Antrasis vykdymas perskaito tai, ką pirmasis paliko duomenų bazėje --- teisės aktus su jų
sprendimais ir teisės aktų nuorodų patikrą. Įmonę jis perskaito iš naujo, tokią, kokia ji yra
antrojo vykdymo pradžioje, ir antrą kartą užfiksuoja profilį (\code{freeze_profile}), kad
ataskaita rodytų tą įmonę, pagal kurią vertintos jos išvados: leidimas, patikrintas ataskaitai
laukiant, ar tuo metu pridėtas dokumentas įskaitomas. Jo \code{documents} etapas vykdomas dar
kartą ir dokumentų inventorius randa podėlyje.
Kai baigiasi paskutinis etapas, ataskaita yra \code{done}; nesėkmė nustato jai būseną
\code{failed} su naudotojui suprantamu pranešimu, atšaukimas --- \code{cancelled}. Baigta
ataskaita skaitoma per \route{GET}{/api/reports/{id}} ir eksportuojama per
\route{GET}{/api/reports/{id}/export}. Šis procesas aptariamas \cref{ch:report}.

\subsection{Programos dokumentas}

Programos dokumentas rengiamas taip pat, \code{programme.generate} rūšies užduotimi
(\cref{fig:arch:product}). \route{POST}{/api/companies/{id}/programme} su rūšimi ---
\code{policy}, \code{governance} arba \code{risks} --- patikrina, ar įmonė turi perskaitomų
dokumentų, ar tos rūšies dokumentas šiuo metu negeneruojamas ir ar liko biudžeto, sukuria
kitą dokumento versiją būsenoje \code{queued} ir įtraukia užduotį į eilę. Pokalbio įrankis
\code{startWorkProduct} kviečia tą patį maršrutą.

\begin{figure}[H]
\centering
\begin{tikzpicture}[x=1mm, y=1mm,
    cell/.style={minimum height=5.6mm, inner ysep=1.5pt, inner xsep=2pt, text width=30mm}]
  \node[svc, text width=17mm] (hc) at (10,0) {konsolė ar pokalbio įrankis};
  \node[svc] (hd) at (48,0) {API (Django)};
  \node[svc] (hw) at (92,0) {vykdytojo ciklas};
  \node[svc, text width=17mm] (ha) at (140,0) {asistento paslauga};
  \draw[draw=rule, line width=.6pt] (10,-6) -- (10,-72);
  \draw[draw=rule, line width=.6pt] (48,-5) -- (48,-72);
  \draw[draw=rule, line width=.6pt] (92,-5) -- (92,-72);
  \draw[draw=rule, line width=.6pt] (140,-6) -- (140,-72);

  \draw[flow] (10,-14) -- node[lbl, above] {\code{POST}~\code{…/programme}} (48,-14);
  \draw[flow] (48,-21) -- node[lbl, above] {užduotis eilėje, tada paimta} (92,-21);
  \node[ai, cell] at (92,-28) {documents};
  \node[ai, cell] (calls) at (92,-35) {du modelio iškvietimai};
  \node[chk, cell] (refs) at (92,-42) {references};
  \draw[call] (calls.east) -- node[lbl, above] {\code{/internal/llm/generate}} (140,-35);
  \draw[opt] (refs.east) -- node[lbl, below] {nuorodų vertintojas} (140,-42);
  \draw[flow] (92,-51) -- node[lbl, above] {dokumentas yra projektas} (48,-51);
  \draw[call] (10,-58) -- node[lbl, above] {tikrina \code{GET}~\code{/api/programme/{id}}} (48,-58);
  \node[man, text width=17mm] (person) at (10,-68) {žmogus redaguoja, patvirtina};
  \draw[flow] (person.east) -- node[lbl, above] {\code{PUT}~\code{…/content}, \code{PATCH}} (48,-68);
\end{tikzpicture}
\actorlegend
\caption{Programos dokumentas. Užduotis perskaito įmonės dokumentus, du kartus klausia modelio
  ir atsakymą patikrina pagal dokumento rūšies taisykles; rezultatas yra projektas. Patvirtinti
  jį gali tik žmogaus užklausa.}\label{fig:arch:product}
\end{figure}

Doroklė (\code{generate} faile \file{programme/run.py}) turi tris pakopas. Etapas
\code{documents} surenka įmonės profilį, dokumentus ir kiekvieno dokumento inventorių --- iš
podėlio arba, jei jo dar nėra, modelio iškvietimu (\code{ensure_inventory}). Du pačios rūšies
etapai atlieka du modelio iškvietimus, antrajam perduodama tai, ką parašė pirmasis. Etapas
\code{references} suteikia atsakymui rūšies formą ir jos taisykles patikrina kodu; tekste
įvardytos nuostatos ieškomos teisės aktų bazėje, o ar kiekviena iš jų reglamentuoja tai, ką
apie ją sako tekstas, vieną kartą įvertina modelis. Tada dokumentas įrašomas kaip
\code{draft} su turiniu ir kiekvienos taisyklės rezultatu. Iš šios būsenos jis niekada
neišeina savaime: žmogus jį redaguoja ir patvirtina per API (\cref{ch:programme}).

\subsection{Teisės aktų bazės atnaujinimas}

Teisės aktų bazė atnaujinama \code{corpus.ingest} rūšies užduotimis, po vieną kiekvienam
teisės aktui (\cref{fig:arch:refresh}). Administratorius tokią užduotį įtraukia į eilę maršrutu
\route{POST}{/api/corpus/documents/{id}/refresh}; kartą per savaitę planuoklis į eilę įtraukia
\code{corpus.refresh_all}, kuri tik įtraukia po vieną \code{corpus.ingest} kiekvienam įjungtam
teisės aktui, kurio tekstas atsisiunčiamas, o ne įkeliamas.

\begin{figure}[H]
\centering
\begin{tikzpicture}[x=1mm, y=1mm,
    cell/.style={minimum height=5.6mm, inner ysep=1.5pt, inner xsep=2pt, text width=34mm}]
  \node[svc, text width=22mm] (hc) at (13,0) {administratorius arba savaitinis tvarkaraštis};
  \node[svc] (hw) at (64,0) {vykdytojo ciklas};
  \node[ext, text width=20mm] (hs) at (108,0) {šaltinis, per išėjimą};
  \node[svc, text width=17mm] (ha) at (140,0) {asistento paslauga};
  \draw[draw=rule, line width=.6pt] (13,-6) -- (13,-64);
  \draw[draw=rule, line width=.6pt] (64,-5) -- (64,-64);
  \draw[draw=rule, line width=.6pt] (108,-6) -- (108,-64);
  \draw[draw=rule, line width=.6pt] (140,-6) -- (140,-64);

  \draw[flow] (13,-14) -- node[lbl, above] {\code{corpus.ingest} eilėje} (64,-14);
  \draw[call] (64,-21) -- node[lbl, above] {atsisiųsti} (108,-21);
  \node[chk, cell] at (64,-28) {nepakito? tada sustoti čia};
  \node[svc, cell] at (64,-35) {išgauti, atpažinti, skaidyti};
  \node[ai, cell] (emb) at (64,-42) {vektorizuoti};
  \node[chk, cell] at (64,-49) {aktyvuoti redakciją};
  \draw[call] (emb.east) -- node[lbl, above, pos=.72] {\code{/internal/llm/embed}} (140,-42);
  \draw[flow] (64,-59) -- node[lbl, above] {teisės aktas turi naują galiojančią redakciją} (13,-59);
\end{tikzpicture}
\actorlegend
\caption{Vieno teisės akto atnaujinimas. Užduočių vykdytojas per išėjimą atsisiunčia šaltinį
  ir sustoja, jei niekas nepasikeitė. Kitu atveju jis šalia galiojančios redakcijos sukuria
  naują ir viena transakcija persijungia į ją.}\label{fig:arch:refresh}
\end{figure}

Doroklė (\code{ingest} faile \file{corpus/ingest.py}) registruoja penkis etapus. \code{fetch}
gauna baitus iš teisės akto šaltinio ir palygina juos su galiojančia redakcija: ta pati turinio
maiša arba ta pati leidimo žyma iškart užbaigia užduotį atsakymu \enquote{nepakito}, todėl
savaitinė nepakitusios teisės aktų bazės patikra nieko nekainuoja. Kitu atveju sukuriama
nauja redakcijos eilutė, o \code{extract} vaikiniame procese su atminties ir laiko ribomis
perskaito leidinio struktūrą, \code{parse} įrašo nuostatas kaip struktūros vienetus,
\code{chunk} suskaido juos paieškai į fragmentus, o \code{embed} fragmentus vektorizuoja. Tik
tada redakcija paverčiama galiojančia kartu su įrašu apie tai, ką ji pakeitė, palyginti su ankstesne,
vienoje transakcijoje; iki tos akimirkos kiekvienas skaitytojas mato visą senąją redakciją.
Nepavykęs vykdymas palieka savo redakciją pažymėtą kaip nepavykusią, o pakartojimas pradeda
naują. Įkelto standarto atveju \code{extract} etapas taip pat gali paprašyti modelio pažymėti
antraštes. Visa tai aptariama \cref{ch:corpus}.

\section{Foninės užduotys}\label{sec:arch:jobs}

Viskas, kas trunka ilgiau nei užklausa, yra \term{užduotis}: viena lentelės eilutė, kurią
paima ir vykdo \term{užduočių vykdytojas}. Nėra nei pranešimų tarpininko, nei atskiro
planuoklio. Eilė yra dvi lentelės toje pačioje duomenų bazėje, kurioje laikoma visa kita, o
kodas yra API programa \code{jobs}: \file{jobs/models.py} (lentelės), \file{jobs/queue.py}
(įtraukimas į eilę), \file{jobs/registry.py} (rūšys), \file{jobs/runner.py} (paėmimas ir
vykdymas), \file{jobs/schedule.py} (periodinės rūšys) ir
\file{jobs/management/commands/run_worker.py} (ciklas).

\subsection{Eilė: dvi lentelės}

Užduotis yra \code{jobs} lentelės eilutė (\cref{tab:arch:jobrow}). Jos \term{rūšis} įvardija
doroklę, kuri ją vykdys, turinys yra doroklės įvestis, o \term{objektas} nurodo, apie ką ji,
--- \code{report:<id>}, \code{corpus:<id>}, \code{document:<id>}, --- kad puslapiai galėtų
išvardyti vieno dalyko užduotis, o maršrutas matytų, kad tokia užduotis jau vykdoma.
\code{enqueue} (\file{jobs/queue.py}) įrašo eilutę; ji atmeta rūšį, kuri neturi doroklės,
kad klaidingai parašyta rūšis sukeltų klaidą ten, kur ji parašyta, ir niekada amžinai
nelauktų lentelėje.

\begin{table}[H]
\caption{\code{jobs} lentelės stulpeliai, paverčiantys ją eile}\label{tab:arch:jobrow}
\begin{tabularx}{\textwidth}{@{}L{37mm}Y@{}}
\toprule
\hd{Stulpelis} & \hd{Reikšmė} \\
\midrule
\code{kind}, \code{payload}, \code{subject} & Kuri doroklė, su kokia įvestimi, apie ką. \\
\code{status} & \code{queued}, \code{running}, \code{done}, \code{failed} arba \code{cancelled}; tikrinimo apribojimas kitokios reikšmės neleidžia. \\
\code{priority}, \code{run_after} & Paėmimo tvarka ir ankstyviausia akimirka, kada užduotį galima paimti. \\
\code{attempts}, \code{max_attempts} & Kiek kartų užduotis buvo paimta ir kiek kartų leidžia jos rūšis. \\
\code{locked_by}, \code{heartbeat_at} & Kuris užduočių vykdytojas laiko užduotį ir kada jis paskutinį kartą pranešė, kad gyvas. \\
\code{cancel_requested} & Nustato API; skaito doroklė. \\
\code{progress} & \texttt{\{step, percent, message\}}: ką puslapiai rodo, kol užduotis vykdoma. \\
\code{result}, \code{error} & Ką grąžino doroklė arba kodėl ji nepavyko. \\
\code{created_by} & Naudotojas, dėl kurio atsirado užduotis; planuoklio užduotims tuščia. Užduoties modelio iškvietimai priskiriami šiam naudotojui, jei apskritai kam nors priskiriami. \\
\code{created_at}, \code{started_at}, \code{finished_at} & Trys jos gyvavimo akimirkos. \\
\bottomrule
\end{tabularx}
\end{table}

Doroklė savo darbą dalija į \term{etapus}, ir kiekvienas etapas yra \code{job_steps} eilutė:
jo eilės numeris ir pavadinimas, būsena (\code{running}, \code{done} arba \code{failed}, kaip
jas įrašo vykdymo modulis), nuoroda į įvestį, išvestis, trukmė ir klaida --- bei jame atliktų
modelio iškvietimų modelis, instrukcijos versija ir sąnaudos. Etapas yra vienetas, kuriuo
skaitoma ataskaitos kaina ir kilmė: \enquote{spragų analizė vykdyta su šia šios instrukcijos
versija, atliko tiek iškvietimų ir kainavo tiek}.

\begin{codebox}{Ką etapas įrašo apie savo modelio iškvietimus}
Doroklė niekada pati neskaičiuoja savo žetonų. Kiekvienas modelio iškvietimas
registruojamas su etapo, kuriam jis atliktas, id, o etapui pasibaigus --- sėkmingai ar ne,
--- vykdymo modulis viena grupuojančia užklausa perskaito etapo eilutes iš \code{llm_calls}
(\code{_usage_from_calls} faile \file{jobs/runner.py}) ir įrašo į etapą: iškvietimų skaičių,
įvesties ir išvesties žetonus ir kainą, iš viso ir pagal modelį; modelį, atlikusį daugiausia
iškvietimų; ir instrukcijos versiją, jei visi iškvietimai vykdyti su viena. Todėl ir pusiaukelėje
nepavykęs etapas rodo, kiek jau buvo išleidęs. Instrukcijos versija laikoma paprastu tekstu, o
ne išoriniu raktu, todėl įrašas išlieka ir ištrynus versiją, kurią jis įvardija.
\end{codebox}

\subsection{Užduoties gyvavimas}

\Cref{fig:arch:jobstates} rodo užduoties būsenas ir kiekvieną perėjimą iš vienos į kitą. Likusi
šio poskyrio dalis juos aptaria.

\begin{figure}[H]
\centering
\begin{tikzpicture}[x=1mm, y=1mm]
  \node[svc, minimum width=22mm] (q) at (14,0) {eilėje};
  \node[svc, minimum width=22mm] (r) at (74,0) {vykdoma};
  \node[svc, minimum width=22mm] (d) at (136,15) {atlikta};
  \node[svc, minimum width=22mm] (f) at (136,0) {nepavyko};
  \node[svc, minimum width=22mm] (c) at (136,-15) {atšaukta};

  \draw[flow] (q) -- node[lbl, above] {paimta: \code{attempts} + 1} (r);
  \draw[flow] (r.north east) -- node[lbl, pos=.55] {doroklė grąžina rezultatą} (d.west);
  \draw[flow] (r.east) -- node[lbl, pos=.5, align=center] {klaida ar žuvęs vykdytojas,\\ bandymų neliko} (f.west);
  \draw[flow] (r.south east) -- node[lbl, pos=.55] {atšaukimas pastebėtas patikroje} (c.west);
  \draw[flow] (r.north) to[out=150, in=30] (q.north);
  \node[note, anchor=south, align=left] at (44,15) {atgal į eilę:\\ klaida, kai liko bandymų (po 30 s)\\ vykdytojas stabdomas (bandymas grąžinamas)\\ gyvybės signalas senesnis nei 120 s, liko bandymų};
  \draw[flow] (q.south) |- node[lbl, pos=.75] {atšaukta dar eilėje} ([yshift=-2mm]c.west);
\end{tikzpicture}
\caption{Užduoties būsenos. Užduotis paimama iš eilės ir baigiasi kaip atlikta, nepavykusi
  arba atšaukta; trys įvykiai vykdomą užduotį vietoj to grąžina į eilę.}\label{fig:arch:jobstates}
\end{figure}

\paragraph{Paėmimas ir užrakinimas.} Užduočių vykdytojas paima po vieną užduotį (\code{claim}
faile \file{jobs/runner.py}): iš eilučių, kurios yra \code{queued} ir kurių \code{run_after}
jau praėjo, tą, kurios \code{priority} didžiausias, o tarp lygių --- seniausią. Eilutė
parenkama PostgreSQL konstrukcija \texttt{FOR UPDATE SKIP LOCKED}, kuri ją užrakina
transakcijos trukmei, o bet kuris kitas tą pačią akimirką klausiantis vykdytojas ją praleidžia
ir paima kitą. Toje pačioje transakcijoje eilutei nustatoma būsena \code{running},
\code{locked_by} nustatomas vykdytojo id (konteinerio serverio vardas ir aštuoni atsitiktiniai
simboliai), \code{heartbeat_at} --- dabarties laikas, o \code{attempts} padidinamas vienetu.
Todėl du vykdytojai niekada negali vykdyti tos pačios užduoties, ir keli greta veikiantys
ciklai yra saugūs. Užklausą aptarnauja indeksas pagal būseną, \code{run_after}, prioritetą ir
sukūrimo laiką.

\paragraph{Prioritetas ir atidėjimas.} Prioritetas yra paprastas sveikasis skaičius,
numatytasis 0, didesnis vykdomas pirmiau. Kode naudojamos keturios reikšmės: 5 --- įkelto
dokumento indeksavimui, kad dokumentas, kurį naudotojas ką tik įmetė į pokalbį, būtų
pasiekiamas paieškai, dar prieš ateinant eilei ilgai teisės aktų bazės užduočiai; 0 --- viskam, ką
naudotojas ar administratorius pradeda per API (ir komandos \code{seed_corpus} įkėlimo
užduotims); $-1$ --- užduotims, kurias masinės komandos \code{extract_obligations},
\code{verify_obligations} ir \code{read_sanctions} įtraukia visai teisės aktų bazei
(\code{PRIORITY} kiekvienoje), kad naudotojo pradėtas darbas eitų pirmas; ir $-5$ ---
planuoklio užduotims. \code{run_after} atideda užduotį: pakartojimas nukeliamas 30~sekundžių
į priekį (\code{BACKOFF_S}), o indeksavimo užduotis, radusi savo naudotojo dienos biudžetą
išnaudotą, vėl įtraukia save į eilę su vykdymo laiku penkios minutės po UTC vidurnakčio, kai biudžetas
atsinaujina.

\paragraph{Bandymai.} Kiekviena rūšis užregistruojama su leidžiamu bandymų skaičiumi
(\cref{tab:arch:kinds}). Kai doroklė iškelia klaidą, o užduočiai liko bandymų, užduotis su
klaidos tekstu ir atidėjimu grįžta į eilę; kitu atveju ji baigiasi kaip \code{failed}.
Dauguma rūšių leidžia vieną bandymą; teisės akto atsisiuntimas (\code{corpus.ingest}) leidžia
du, o dokumento indeksavimas (\code{documents.index}) --- tris. Klaida, iškelta, kai užduoties
atšaukimo žymė nustatyta, užduotį vietoj to užbaigia kaip \code{cancelled}, kad ir kiek
bandymų liktų.

\paragraph{Gyvybės signalas ir žuvęs užduočių vykdytojas.} Kol doroklė vykdoma, antra gija
kas 15~sekundžių pažymi \code{heartbeat_at} (\code{HEARTBEAT_S}). Kiekvieno savo ciklo
apsisukimo pradžioje vykdytojas ieško vykdomų užduočių, kurių gyvybės signalas senesnis nei
120~sekundžių (\code{STALE_S}): jų vykdytojas laikomas žuvusiu (\code{recover_stale}). Tokia
užduotis grįžta į eilę, jei jai liko bandymų. Jei neliko, jai nustatoma būsena \code{failed} su
pranešimu, kad jos vykdytojas netikėtai sustojo, kartu nepavykusiais pažymimi ir jos dar
atviri etapai --- kiekvienas su tuo, ką buvo sunaudoję jo modelio iškvietimai, --- ir
paleidžiamas rūšies kablys (angl. \en{hook}) \code{on_abandon}. Kablys skirtas dviem rūšims,
kurių užduotį atspindi atskira jos pačios eilutė: jis ataskaitai arba programos dokumentui taip
pat nustato nepavykusią būseną, kad nė vienas neliktų amžinai \enquote{vykdomas} be jokios
užduoties už jo.

\begin{codebox}{Aptvėrimas: jokio įrašo be užrakto}
Užduočių vykdytojas, kuris tik lėtai veikė --- buvo sustojęs kelioms minutėms, bet nežuvo, ---
gali pabusti, kai kitas vykdytojas jau atkūrė jo užduotį. Todėl kiekvienas įrašas, kuris
užbaigia užduotį --- atlikta, nepavyko, atšaukta, atgal į eilę, --- vykdomas tik su sąlyga,
kad eilutė vis dar \code{running} ir vis dar turi to vykdytojo \code{locked_by}
(\code{_fenced} faile \file{jobs/runner.py}), o eigos įrašas --- su sąlyga, kad eilutė vis dar
turi tą \code{locked_by}; atnaujinimas, neatitinkantis jokios eilutės, nieko nekeičia. Gyvybės
signalo gija tai pastebi pirmoji: kai jos atnaujinimas neatitinka jokios eilutės, ji nustato žymę,
doroklė per kitą patikrą sustoja su \code{JobLost}, ir šiam bandymui daugiau niekas
neįrašoma.
\end{codebox}

\paragraph{Atšaukimas.} \route{POST}{/api/jobs/{id}/cancel} (\file{jobs/api/views.py}) daro
vieną iš dviejų dalykų. Dar eilėje esančiai užduočiai iškart nustatoma būsena
\code{cancelled}, o jos rūšies kablys \code{on_abandon} sutvarko už jos esančią ataskaitą ar
programos dokumentą. Vykdomai užduočiai tik nustatoma žymė \code{cancel_requested}; doroklė
ją pamato kito \code{ctx.check_cancelled()} kvietimo metu (doroklės jį atlieka tarp brangių
savo fazių), iškelia \code{JobCancelled}, ir užduotis baigiasi kaip \code{cancelled} --- ne
kaip nesėkmė ir be pakartojimo. Abu įrašai yra sąlyginiai ir atliekami šia tvarka, todėl
užduotis, kurią vykdytojas paima tarp jų, vis tiek gauna žymę. Jau pasibaigusi užduotis atsako
409. Ataskaitos ir programos dokumentai turi savus atšaukimo maršrutus, kurie su savo užduotimi
daro tą patį. Naudotojas gali atšaukti savo sukurtas užduotis, administratorius --- bet kurią;
ta pati taisyklė lemia, kurias užduotis išvardija \route{GET}{/api/jobs}.

\paragraph{Eiga.} Doroklė praneša, kur yra, kvietimu \texttt{ctx.progress(step, percent,
message)}, kuris įrašo tris reikšmes į užduoties eilutę. Užduočių puslapis bei ataskaitos ir
programos dokumento eigos skydeliai periodiškai tikrina užduotį ir jas rodo.

\subsection{Užduočių rūšys}

Doroklė yra įprasta funkcija, užregistruota savo rūšies vardu dekoratoriumi
\code{@register} (\file{jobs/registry.py}); \file{jobs/handlers.py} importuoja kiekvieną
modulį, kuris ją registruoja, ir žiniatinklio procese, ir užduočių vykdytojuje, nes
žiniatinklio procesas turi žinoti rūšis, kad galėtų jas įtraukti į eilę. \Cref{tab:arch:kinds}
išvardija visas.

\begin{xltabular}{\textwidth}{@{}L{37mm}YL{43mm}@{}}
\caption{Užregistruotos užduočių rūšys}\label{tab:arch:kinds}\\
\toprule
\hd{Rūšis, doroklė, bandymai} & \hd{Ką daro} & \hd{Kas gali pradėti} \\
\midrule
\endfirsthead
\toprule
\hd{Rūšis, doroklė, bandymai} & \hd{Ką daro} & \hd{Kas gali pradėti} \\
\midrule
\endhead
\bottomrule
\endlastfoot
\code{report.analyze}\newline \file{analysis/pipeline.py}\newline 1 bandymas, atsisakymo kablys & Vienas ataskaitos vykdymas: pirmoji jos dalis, antroji arba abi (\cref{fig:arch:report}). & Bet kuris naudotojas, įmonei, kurią gali matyti: \route{POST}{/api/reports} ir apimties patvirtinimas; tą patį daro pokalbio įrankiai. \\
\code{programme.generate}\newline \file{programme/run.py}\newline 1 bandymas, atsisakymo kablys & Nauja vieno programos dokumento versija. & Bet kuris naudotojas, įmonei, kurią gali matyti: \route{POST}{/api/companies/{id}/programme}; pokalbio įrankis. \\
\code{documents.index}\newline \file{companies/indexing.py}\newline 3 bandymai & Suskaido vieną įkeltą dokumentą į ištraukas ir jas vektorizuoja. Jei įkėlusiojo biudžetas išnaudotas, laukia kitos dienos. & Bet kuris naudotojas, įkeldamas dokumentą (\route{POST}{/api/documents}); pati užduotis, kai laukia. \\
\code{corpus.ingest}\newline \file{corpus/ingest.py}\newline 2 bandymai & Atsisiunčia vieną teisės aktą arba perskaito įkeltą failą ir paverčia jį galiojančia akto redakcija (\cref{fig:arch:refresh}). & Administratorius: pridėdamas teisės aktą su šaltiniu, jį atnaujindamas, įkeldamas failą. Žemiau esanti rūšis. Komanda \code{seed_corpus}. \\
\code{corpus.refresh_all}\newline \file{corpus/ingest.py}\newline 1 bandymas & Įtraukia po vieną \code{corpus.ingest} kiekvienam įjungtam atsisiunčiamam teisės aktui. & Planuoklis, kas savaitę. Administratorius: \route{POST}{/api/corpus/refresh-all}. \\
\code{corpus.reembed}\newline \file{corpus/ingest.py}\newline 1 bandymas & Iš naujo vektorizuoja kiekvieną galiojančių redakcijų fragmentą ir, jei nenurodytas vienas teisės aktas, kiekvieną pareigą abiem kalbomis ir kiekvieną įmonės dokumentą. Reikalinga pasikeitus vektorizavimo modeliui. & Administratorius: \route{POST}{/api/corpus/reembed}. \\
\code{corpus.extract_obligations}\newline \file{corpus/obligations.py}\newline 1 bandymas & Išskiria pareigas iš vieno teisės akto galiojančios redakcijos. & Administratorius, kiekvienam aktui atskirai. Komanda \code{extract_obligations} --- kiekvienam atsilikusiam aktui. \\
\code{corpus.verify_obligations}\newline \file{corpus/verify.py}\newline 1 bandymas & Paveda vertintojui palyginti vieno teisės akto pareigas su nuostatomis, kurias jos cituoja. & Administratorius, kiekvienam aktui atskirai. Komanda \code{verify_obligations}. \\
\code{corpus.read_sanctions}\newline \file{corpus/sanctions.py}\newline 1 bandymas & Perskaito dar neperskaitytų teisės aktų sankcijų nuostatas, vieną aktą po kito. & Tik komanda \code{read_sanctions}. \\
\code{ops.maintenance}\newline \file{ops/maintenance.py}\newline 1 bandymas & Kasdieniai valymai (\cref{tab:arch:purges}). & Planuoklis, kasdien. \\
\code{ops.noop}\newline \file{jobs/handlers.py}\newline 1 bandymas & Dūmų testas: keli tušti etapai su pasirinktine pauze ir priverstine nesėkme. & Testų rinkinys; žmogus komandinėje eilutėje. \\
\end{xltabular}

Joks API maršrutas neįtraukia į eilę bet kokios rūšies užduoties. Užduotis atsiranda tik kaip
šalutinis maršruto, patikrinusio kviečiantįjį ir objektą, poveikis, kaip serveryje paleistos
valdymo komandos rezultatas arba kaip planuoklio veiksmas. Maršrutai, kurie atnaujina teisės aktą,
išskiria jo pareigas ar paveda vertintojui jas perskaityti, užuot pridėję antrą užduotį,
grąžina tam aktui jau eilėje esančią ar vykdomą užduotį.

Ne kiekviena ilga operacija lentelėje turi savo rūšį: tai, ko reikia ataskaitai, atliekama
pačios ataskaitos užduotyje. Ataskaita, kurios teisės aktuose pareigos dar niekada nebuvo
išskirtos, jas išskiria savo \code{obligations} etape ir ten pat perskaito jų sankcijų
nuostatas; ji nelaukia kitos užduoties.

\subsection{Užduočių vykdytojas ir jo ciklas}

Ciklas yra API valdymo komanda \code{run_worker}. Kiekvieno apsisukimo metu jis pažymi gyvybės
signalo failą, kartą per minutę leidžia pažiūrėti planuokliui, atkuria užduotis, kurių
vykdytojas žuvo, paima vieną užduotį ir vykdo ją iki galo; jei nėra ko paimti, laukia dvi
sekundes (\code{--poll}) ir žiūri vėl. Vienas ciklas vienu metu vykdo vieną užduotį.

\paragraph{Kodėl užduočių vykdytojas yra Docker klientas.} Ciklui reikia lygiai to, ką turi
API: to paties kodo, tų pačių parametrų ir paslapčių, to paties įkeltų failų katalogo, to
paties ryšio su duomenų baze, tų pačių susietų adresų teisės aktams atsisiųsti. Compose failas
visa tai suteikia, paleisdamas ciklą API konteinerio \emph{viduje}. Konteineris
\code{legalops-worker} yra tik standartinis Docker komandinės eilutės atvaizdas su viena
fiksuota komanda \texttt{docker exec legalops-django python3 manage.py run\_worker} ir į jį
prijungtu serverio Docker lizdu. Jis visai neturi tinklo, o jo žurnalas yra ciklo žurnalas.
Niekas nekuriama ir nekonfigūruojama dukart, ir API kodo ar aplinkos pakeitimas negali
palikti užduočių vykdytojo atsilikusio. Kaina --- lizdas: prieiga prie jo prilygsta root
teisėms serveryje, todėl tas konteineris neturi tinklo ir nieko kito nevykdo
(\cref{ch:security}).

\paragraph{Veikiančio ciklo perėmimas.} Ši tvarka turi vieną nepatogią savybę, kurią kodas
tvarko aiškiai. Ciklas, paleistas \texttt{docker exec}, priklauso API konteineriui, o ne
užduočių vykdytojo konteineriui. Kai API konteineris paleidžiamas iš naujo, ciklas žūva kartu
su juo, \texttt{docker exec} grįžta, vykdytojo konteineris baigia darbą, o jo paleidimo iš
naujo politika jį --- ir kartu naują ciklą --- vėl paleidžia. Tačiau kai iš naujo paleidžiamas
tik vykdytojo konteineris, senasis ciklas API konteineryje toliau veikia, o šalia jo tuoj bus
paleistas antras. Todėl naujas ciklas pirmiausia perima darbą (\code{take_over} faile
\file{jobs/management/commands/run_worker.py}): jis
perskaito proceso id, kurį ankstesnis ciklas paliko faile (\file{/tmp/worker.pid}),
\file{/proc} patikrina, ar šis procesas tikrai yra \code{run_worker} ciklas, o ne nesusijęs
procesas, atsitiktinai turintis tą patį numerį, siunčia jam sustabdymo signalą ir įrašo į
failą savo id. Jis nelaukia: senasis ciklas iškart nustoja imti užduotis ir užbaigia tik tą
etapą, kuriame yra.

\paragraph{Tvarkingas sustabdymas.} Sustabdymo signalas tik nustato žymę. Ciklas nustoja imti
užduotis, o vykdoma doroklė žymę pamato per kitą patikrą ir iškelia \code{WorkerStopping}.
Tada vykdymo modulis grąžina užduotį į eilę ir grąžina jos bandymą (\code{attempts}
sumažinamas vienetu), kad pertraukimas nebūtų laikomas nesėkme; kitas ciklas užduotį vykdo
iš naujo. Ataskaitai tai reiškia, kad pertraukta dalis pradedama iš naujo: doroklė pradeda
nuo eilučių, kurias paliko ankstesnis tos pačios dalies bandymas, ištrynimo, todėl niekas
nepasidvigubina.

\begin{notebox}{Ką sistemos atnaujinimas daro vykdomai užduočiai}
Du būdai, kuriais ciklas gali baigtis, turi skirtingas pasekmes, ir skirtumas svarbus
rūšims, kurios leidžia vienintelį bandymą. Ciklas, kurio \emph{paprašoma} sustoti --- nes
darbą perima naujas ciklas, --- grąžina savo užduotį nepanaudojęs bandymo, ir užduotis
vykdoma vėl. Ciklas, kuris \emph{priverstinai nutraukiamas} kartu su API konteineriu, kaip nutinka, kai visa
sistema sustabdoma atnaujinimui, nieko grąžinti negali: jo užduotis lieka \code{running}, kol
jos gyvybės signalui sukanka dvi minutės, ir tada ją atkuria kitas ciklas.
\code{corpus.ingest} ir \code{documents.index} turi likusių bandymų ir vėl įtraukiamos į
eilę. \code{report.analyze} ar \code{programme.generate} užduotis savo vienintelį bandymą
išnaudojo, todėl baigiasi kaip \code{failed} su pranešimu, kad jos vykdytojas netikėtai
sustojo, o jos ataskaitai ar programos dokumentui kartu nustatoma nepavykusi būsena, užuot
palikus juos kaboti. Tai, už ką pertrauktas vykdymas jau sumokėjo, išlieka jo etapuose ir
\code{llm_calls}; pačią ataskaitą reikia pradėti iš naujo.
\end{notebox}

\subsection{Planuoklis}

Dvi rūšys yra periodinės, ir \file{jobs/schedule.py} visą tvarkaraštį laiko kaip dviejų
įrašų lentelę: \code{ops.maintenance} kasdien ir \code{corpus.refresh_all} kas septynias
dienas. Kartą per minutę (\code{SCHEDULE_EVERY_S}) ciklas kviečia \code{tick}, kuri suranda
naujausią kiekvienos rūšies užduotį ir įtraukia į eilę naują, su prioritetu $-5$ ir objektu
\code{schedule:<kind>}, nebent ta naujausia užduotis dar eilėje ar vykdoma arba buvo sukurta
per intervalą. Taigi tvarkaraštis saugomas pačioje eilėje, kurią jis pildo. Nėra cron
konteinerio; planuoklis niekada neįtraukia rūšies, kol viena iš jos užduočių dar laukia ar
vykdoma; administratoriaus rankiniu būdu pradėtas atnaujinimas laikomas tos savaitės vykdymu;
o po ilgo sutrikimo kiekvienos rūšies laikas tiesiog būna atėjęs, ir ji įvykdoma vieną kartą.
Šio patikrinimo nesėkmė užregistruojama žurnale ir niekada nesustabdo ciklo.

\subsection{Kasdienė priežiūra}

\code{ops.maintenance} (\file{ops/maintenance.py}) pašalina tai, kas nebeturi jokios
naudos, ir atsako, kiek ko pašalino (\cref{tab:arch:purges}).

\begin{table}[H]
\caption{Ką pašalina kasdienė priežiūros užduotis}\label{tab:arch:purges}
\begin{tabularx}{\textwidth}{@{}L{38mm}YL{43mm}@{}}
\toprule
\hd{Kas} & \hd{Sąlyga} & \hd{Konstanta} \\
\midrule
Sesijos & Pasibaigusios. & --- \\
Baigtos užduotys su jų etapais & Atliktos, nepavykusios ar atšauktos prieš daugiau nei 90 dienų --- išskyrus užduotį, į kurią rodo ataskaita: ji saugoma tol, kol yra ataskaita. & \code{JOB_RETENTION_DAYS}~=~90 \\
DI sąnaudų įrašai & \code{llm_calls} eilutės, senesnės nei 400 dienų, kartu su jose esančiomis pokalbių stenogramomis. & \code{TELEMETRY_RETENTION_DAYS}~=~400 \\
Niekam nepriskirti įkelti failai & Dokumentų katalogo failai, kurių neįvardija jokia dokumento eilutė ir kurie senesni nei šešios valandos: tai, ką palieka gedimas tarp failo ir eilutės įrašymo. & \code{ORPHAN_GRACE_HOURS}~=~6 \\
Užklausų ribojimo langai & Praėjęs jų nunulinimo laikas. & --- \\
Nepanaudoti pokalbio įkėlimai & Dokumentai, nepriskirti jokiai įmonei, įkelti prieš daugiau nei septynias dienas ir neįvardyti jokioje pokalbio žinutėje: failas, įmestas į pokalbį, bet taip ir neišsiųstas. Ne daugiau kaip 500 per vykdymą. & \code{ORPHAN_UPLOAD_DAYS}~=~7, \code{ORPHAN_UPLOAD_BATCH}~=~500 \\
\bottomrule
\end{tabularx}
\end{table}

Antroji eilutė paaiškina, kodėl ataskaita išsaugo savo kilmę: užduotis, į kurią rodo
ataskaita, su etapais, kuriuose įrašyta kiekviena instrukcijos versija, modelis ir kaina,
nevaloma, kol ataskaita egzistuoja.
Ataskaita, parengta su apimties peržiūra, rodo į savo antrosios dalies užduotį. Pirmosios
dalies užduotis išvaloma kaip bet kuri kita, ir po vykdymo jos niekas nebeskaito: ataskaitos
statistika su abiejų dalių instrukcijų versijomis, modeliais ir kaina įrašoma į ataskaitą, kai
vykdymas baigiasi (\cref{ch:data}), o DI sąnaudų įrašai iškvietimus saugo 400 dienų.

\section{Technologijų rinkinys ir du režimai}\label{sec:arch:stack}

\Cref{tab:arch:stack} pateikia saugykloje fiksuotas versijas. Python paketai tiksliai
fiksuoti \file{django/Dockerfile} faile --- projektas neturi reikalavimų failo, --- o baziniai
atvaizdai fiksuoti žyme. Du Node projektai savo \file{package.json} deklaruoja minimalias
versijas; tiksliai įdiegtos versijos yra nurodytos kiekviename \file{package-lock.json}.

\begin{xltabular}{\textwidth}{@{}L{29mm}Y@{}}
\caption{Technologijų rinkinys su saugykloje deklaruotomis versijomis}\label{tab:arch:stack}\\
\toprule
\hd{Dalis} & \hd{Komponentai ir versijos} \\
\midrule
\endfirsthead
\toprule
\hd{Dalis} & \hd{Komponentai ir versijos} \\
\midrule
\endhead
\bottomrule
\endlastfoot
Įėjimas ir išėjimas & Caddy 2.11 (\code{caddy:2.11-alpine}); Caddy su ketvirtojo lygmens moduliu (\code{vuknf/caddy-l4:2.11}). \\
API ir užduočių vykdytojo ciklas & Python 3.13 (\code{python:3.13-slim}); Django 5.2.17; psycopg 3.3.6 ir \code{pgvector} sąsaja 0.5.0; gunicorn 26.2.0; django-environ 0.14.0; bcrypt 5.0.0; requests 2.34.2. Dokumentų skaitymas: pypdf 6.19.0, cryptography 50.0.2, python-docx 1.2.0, lxml 6.1.3, PyYAML 6.0.3. Eksportas: WeasyPrint 70.0, openpyxl 3.1.5, Markdown 3.11. \\
Asistento paslauga & Node 24.16.0; Express 5.1; DI SDK \code{ai} 7.0 su tiekėjų paketais \code{@ai-sdk/openai} 4.0, \code{@ai-sdk/openai-compatible} 3.0 ir \code{@ai-sdk/anthropic} 4.0, taip pat \code{@ai-sdk/mcp} 2.0; zod 4.1. \\
Konsolė ir viešoji svetainė & React 19.2; Vite 7.3; MUI 7.3 su duomenų tinkleliu 8.27; Tailwind CSS 4.2; assistant-ui 0.15; TanStack Query 5; React Router 7.13; axios 1.13. Testuojama su vitest 5 ir msw 2; gamybos režimu sukurtą programą aptarnauja Caddy 2.11. \\
Duomenų bazė & PostgreSQL 17.11 (\code{postgres:17.11-alpine3.24}) su į atvaizdą sukompiliuotu \code{pgvector} 0.8.6. \\
Užduočių vykdytojo konteineris & Docker komandinės eilutės klientas (\code{docker:27-cli}). \\
Duomenų bazės įrankis & DbGate 7.3.3. \\
\end{xltabular}

Compose failas parašytas taip, kad veiktų vienu iš dviejų režimų, o besiskiriančios eilutės
turi komentarą \texttt{\# Dev} arba \texttt{\# Prod}. Norint perjungti režimą, vienos žymės
eilutes reikia užkomentuoti, kitos --- atkomentuoti ir paleisti atnaujinimo scenarijų
(\cref{tab:arch:modes}). Saugykloje esantis failas nustatytas kūrimo režimui.

\begin{table}[H]
\caption{Du \file{docker-compose.yml} režimai}\label{tab:arch:modes}
\begin{tabularx}{\textwidth}{@{}L{27mm}YY@{}}
\toprule
\hd{Paslauga} & \hd{Kūrimo režimas} & \hd{Gamybos režimas} \\
\midrule
Konsolė ir viešoji svetainė & \file{vite/Dockerfile.dev}: Vite kūrimo serveris, veikiantis su iš serverio prijungtu išeities kodo katalogu; išsaugojus failą, puslapis perkraunamas. & \file{vite/Dockerfile}: programa vieną kartą sukuriama iš užfiksuotų priklausomybių, o Caddy rezultatą aptarnauja kaip statinius failus. \\
Asistento paslauga & \file{assistant/Dockerfile.dev}: išeities kodas prijungtas iš serverio; jam pasikeitus, paslaugą iš naujo paleidžia nodemon. & \file{assistant/Dockerfile}: užfiksuotos gamybos priklausomybės ir išeities kodas įdėti į atvaizdą; \texttt{node server.js}. \\
API & Nustatytas \code{DJANGO_RUNSERVER}: Django kūrimo serveris su iš serverio prijungtu išeities kodu. & Kintamasis nenustatytas: gunicorn iš sukurto atvaizdo, numatytuoju atveju trys procesai po aštuonias gijas (\file{django/gunicorn.conf.py}). \\
Tinklai & Konsolės ir asistento paslaugos konteineriai taip pat jungiasi prie išorinio tinklo, kad galėtų įdiegti paketus. & Abu yra tik izoliuotame tinkle. \\
\bottomrule
\end{tabularx}
\end{table}

Dvi kūrimo režimo pasekmės priklauso architektūros skyriui, nes jos susilpnina dvi pagrindines
jo savybes. Konteineris išoriniame tinkle gali tiesiogiai pasiekti internetą, todėl kūrimo
režimu asistento paslaugos neriboja išėjimo leidžiamų adresų sąrašas. Be to, užduočių
vykdytojo ciklas neperkraunamas, kai pasikeičia išeities kodo failas: jis vykdo kodą, su
kuriuo buvo paleistas, kol iš naujo paleidžiamas vykdytojo konteineris, o tada darbą perima
naujas ciklas. Gamybos režimu kodo pakeitimas į bet kurį konteinerį patenka tik per
atnaujinimo scenarijų (\cref{ch:ops}).

\section{Architektūriniai sprendimai}\label{sec:arch:decisions}

Saugykloje užfiksuoti šeši sprendimai, kurie suformavo architektūrą
(\file{docs/markdown/en/architecture.md}, 10~skirsnis). Čia jie pateikiami su priežastimis.

\begin{description}
  \item[API --- Django, o ne FastAPI.] Įraše tai nurodyta kaip pačioje pradžioje priimtas
    sprendimas, be jokio techninio argumento. Kodas rodo, kaip mažai Django naudojama: ORM,
    migracijos, užklausos ciklas ir valdymo komandos, be jokių kartu teikiamų programų.

  \item[Pirmiausia pokalbis; puslapiai peržiūrai ir administravimui.] Numatyta, kad
    naudotojas dirbs kalbėdamasis: aprašys įmonę, įmes jos dokumentus, paprašys analizės.
    Konsolės puslapiai skirti peržiūrėti, kas parengta, ir administruoti sistemą. Pokalbio
    įrankiai apima darbą su įmone ir nieko daugiau --- jie gali skaityti teisės aktų bazę,
    bet ne ją keisti, ir negali pasiekti paskyrų, instrukcijų ar MCP nustatymų. Tai riboja,
    ką dokumente paslėptas nurodymas galėtų priversti pokalbį padaryti (\cref{ch:security}).

  \item[Sistema rengia projektus ir tikrina; tvirtina žmogus.] Ataskaitos apimtis, registro
    pareiga, peržiūrėtojo sprendimas dėl išvados, programos dokumentas: kiekvieną jų
    patvirtina įvardytas asmuo, tai užrašoma veiksmų žurnale, ir sistema nieko nepažymi kaip
    patvirtinto. Balus, lygius ir identifikatorius skaičiuoja kodas, jie niekada neimami iš
    modelio. Sustojimas apimties peržiūrai (\cref{fig:arch:report}) ir projekto būsena
    (\cref{fig:arch:product}) yra šio sprendimo išraiška architektūroje; \cref{ch:control}
    surenka visus tokius taškus.

  \item[Pareigos priklauso teisės aktų bazei ir yra peržiūrimos.] Teisės akto pareigos
    išskiriamos vieną kartą kiekvienai redakcijai ir pakartotinai naudojamos kiekvienoje
    vėlesnėje kiekvienos įmonės ataskaitoje, užuot jas iš naujo skaičius kiekvienai
    ataskaitai. Dėl to ataskaita pigesnė, dvi ataskaitos palyginamos, o pareigos peržiūra turi
    ilgalaikį poveikį. Ataskaita įrašo, kiek jos išvadų remiasi peržiūrėtomis pareigomis. Kaip
    registras buvo sukurtas ir peržiūrėtas, aptariama \cref{ch:obligations}.

  \item[Oficialūs mašininiai teisės šaltiniai.] Tekstai gaunami iš CELLAR ir e-TAR, o ne iš
    EUR-Lex ir e-seimo --- svetainių, kurias skaitytų žmogus, --- nes šios dvi atmeta
    automatinius klientus. \cref{tab:arch:egress} pateiktas išeinančio srauto leidžiamų
    adresų sąrašas yra architektūrinis šio sprendimo pėdsakas.

  \item[Vienas modelių tiekėjas: fakulteto DI tinklų sąsaja.] Kiekvienas modelio iškvietimas
    eina į DI tinklų sąsają adresu \code{ai.knf.vu.lt}. Įraše aiškiai pasakyta, kas dėl to
    palieka serverį: įmonės dokumentai siunčiami į DI tinklų sąsają vektorizuoti ir
    ataskaitos etapuose, o visa, ką perskaito pokalbis, keliauja kartu su jo ėjimais. Visa
    kita --- duomenų bazė, failai, struktūros atpažinimas, patikros, eksportas --- vyksta
    serveryje. Asistento paslaugą galima sukonfigūruoti kitiems tiekėjams, bet tam reikėtų
    atverti jiems leidžiamų adresų sąrašą (\cref{ch:models}).
\end{description}
